% =====================================================================
%  PREAMBLE — biên dịch bằng XeLaTeX (font hệ thống Times New Roman)
% =====================================================================

% ---------------------------------------------------------------- cỡ chữ 13pt
% lớp report chỉ có 10/11/12pt; scrextend (KOMA-Script) đặt lại cỡ chữ gốc thành 13pt
\usepackage[fontsize=13pt]{scrextend}

% ---------------------------------------------------------------- trang
\usepackage[a4paper, left=3cm, right=2cm, top=2.5cm, bottom=2.5cm,
            headheight=15pt, headsep=0.8cm, footskip=1.2cm]{geometry}

% ---------------------------------------------------------------- font & ngôn ngữ
\usepackage{amsmath, amssymb, mathtools}
\usepackage{fontspec}
\usepackage{unicode-math}
\setmainfont{Times New Roman}
\setsansfont{Arial}
\setmonofont{Consolas}[Scale=0.88]
\setmathfont{TeX Gyre Termes Math}
\usepackage{polyglossia}
\setdefaultlanguage{vietnamese}
\setotherlanguage{english}
\PolyglossiaSetup{vietnamese}{indentfirst=true}  % thụt đầu dòng cả đoạn đầu sau tiêu đề

\usepackage{setspace}
\setstretch{1.2}
\usepackage{indentfirst}
\setlength{\parindent}{1cm}
\setlength{\parskip}{3pt plus 1pt}
\usepackage{microtype}
\setlength{\emergencystretch}{3em}  % tránh tràn lề khi dòng có mã nguồn không ngắt được
\usepackage{adjustbox}

% ---------------------------------------------------------------- màu, hình, bảng
\usepackage[table]{xcolor}
\definecolor{ctublue}{HTML}{1F4E8C}
\definecolor{codebg}{HTML}{F5F7FA}
\definecolor{codeframe}{HTML}{C9D3E0}
\definecolor{codekw}{HTML}{1F4E8C}
\definecolor{codestr}{HTML}{A3361B}
\definecolor{codecom}{HTML}{5B6B7B}
\definecolor{tblhead}{HTML}{DCE6F2}
\definecolor{tblrow}{HTML}{F4F7FB}
\definecolor{okgreen}{HTML}{15803D}
\definecolor{badred}{HTML}{B91C1C}

\usepackage{graphicx}
\graphicspath{{figures/}{figures/common/}{figures/results/}{figures/generated/}{figures/ui/}{figures/paper/}}
\usepackage{float}
\usepackage{caption}
\captionsetup{font=small, labelsep=colon, justification=centering, skip=6pt}
\captionsetup[table]{position=top}
\DeclareCaptionListFormat{hinh}{\figurename~#1#2:}
\DeclareCaptionListFormat{bang}{\tablename~#1#2:}
\captionsetup[figure]{listformat=hinh}
\captionsetup[table]{listformat=bang}
\usepackage{subcaption}
\usepackage{booktabs, array, tabularx, longtable}
\renewcommand{\arraystretch}{1.18}
\newcolumntype{L}{>{\raggedright\arraybackslash}X}
\newcolumntype{C}{>{\centering\arraybackslash}X}
\newcolumntype{R}{>{\raggedleft\arraybackslash}X}
\newcolumntype{P}[1]{>{\raggedright\arraybackslash}p{#1}}
\usepackage{siunitx}
\sisetup{group-separator={\,}, group-minimum-digits=5, output-decimal-marker={.},
         detect-all}
\usepackage{enumitem}
\setlist{itemsep=1pt, topsep=3pt, parsep=1pt}
\setlist[itemize,1]{label=--, leftmargin=1.2cm}
\setlist[itemize,2]{label=+, leftmargin=0.8cm}
\setlist[enumerate,1]{leftmargin=1.2cm}

% ---------------------------------------------------------------- TikZ
\usepackage{tikz}
\usetikzlibrary{arrows.meta, positioning, calc, fit, backgrounds, shapes.geometric,
                shapes.misc, matrix, decorations.pathreplacing}
\tikzset{
  >={Stealth[length=2.2mm]},
  box/.style={draw, rounded corners=2pt, align=center, minimum height=8mm,
              inner sep=4pt, fill=white, font=\small},
  pbox/.style={box, fill=tblhead},
  dbox/.style={box, fill=codebg, dashed},
  xorn/.style={draw, circle, inner sep=0pt, minimum size=6mm, font=\small},
  lbl/.style={font=\footnotesize, align=center},
  arr/.style={->, thick},
}
\usepackage{bytefield}

% ---------------------------------------------------------------- mã nguồn
\usepackage{listings}
\lstdefinestyle{py}{
  language=Python,
  basicstyle=\ttfamily\footnotesize\setstretch{1.0},
  keywordstyle=\color{codekw}\bfseries,
  stringstyle=\color{codestr},
  commentstyle=\color{codecom}\itshape,
  numberstyle=\tiny\color{codecom},
  numbers=left, numbersep=6pt, xleftmargin=1.4em, framexleftmargin=1.2em,
  backgroundcolor=\color{codebg}, frame=single, rulecolor=\color{codeframe},
  breaklines=true, columns=fullflexible, keepspaces=true, showstringspaces=false,
  tabsize=4, aboveskip=6pt, belowskip=4pt, upquote=true,
  morekeywords={self, True, False, None, as, with},
}
\lstdefinestyle{plain}{
  basicstyle=\ttfamily\footnotesize\setstretch{1.0}, numbers=none, xleftmargin=0pt,
  framexleftmargin=0pt, commentstyle=\color{codecom}\itshape, morecomment=[l]{\#},
  backgroundcolor=\color{codebg}, frame=single, rulecolor=\color{codeframe},
  columns=fullflexible, keepspaces=true, breaklines=true, aboveskip=6pt, belowskip=4pt,
}
\lstset{style=py}

% ---------------------------------------------------------------- tiêu đề trang
\usepackage{fancyhdr}
\fancypagestyle{noidung}{%
  \fancyhf{}%
  \fancyhead[L]{\small\itshape \TenNgan}%
  \fancyhead[R]{\small\itshape \MaHocPhan}%
  \fancyfoot[C]{\small\thepage}%
  \renewcommand{\headrulewidth}{0.4pt}%
  \renewcommand{\footrulewidth}{0pt}}
\fancypagestyle{dau}{%
  \fancyhf{}%
  \fancyfoot[C]{\small\thepage}%
  \renewcommand{\headrulewidth}{0pt}}

% ---------------------------------------------------------------- tiêu đề chương / mục
\makeatletter
\renewcommand\@makechapterhead[1]{%
  \vspace*{-6pt}%
  {\centering\normalfont\fontsize{16}{21}\selectfont\bfseries
   \ifnum\c@secnumdepth>\m@ne
     \MakeUppercase{\@chapapp}\ \thechapter:\ %
   \fi
   \MakeUppercase{#1}\par}%
  \vspace{16pt}}
\renewcommand\@makeschapterhead[1]{%
  \vspace*{-6pt}%
  {\centering\normalfont\fontsize{16}{21}\selectfont\bfseries\MakeUppercase{#1}\par}%
  \vspace{16pt}}
\renewcommand\section{\@startsection{section}{1}{\z@}%
  {-2.6ex \@plus -0.8ex \@minus -.2ex}{1.0ex \@plus .2ex}%
  {\normalfont\fontsize{14}{17}\selectfont\bfseries}}
\renewcommand\subsection{\@startsection{subsection}{2}{0.6em}%
  {-2.2ex \@plus -0.6ex \@minus -.2ex}{0.8ex \@plus .2ex}%
  {\normalfont\normalsize\bfseries\itshape}}
\renewcommand\subsubsection{\@startsection{subsubsection}{3}{1.2em}%
  {-1.8ex \@plus -0.5ex \@minus -.2ex}{0.6ex \@plus .2ex}%
  {\normalfont\normalsize\itshape}}
\setcounter{secnumdepth}{3}
\setcounter{tocdepth}{2}
\renewcommand\thesubsubsection{\alph{subsubsection}.}
% đánh số hình/bảng/phương trình theo chương
\renewcommand\theequation{\thechapter.\arabic{equation}}
% mục lục: dòng cấp chương in hoa giống tiêu đề chương ("CHƯƠNG 1: ...", "PHỤ LỤC A: ...")
\def\@chapter[#1]#2{%
  \ifnum\c@secnumdepth>\m@ne
    \refstepcounter{chapter}%
    \typeout{\@chapapp\space\thechapter.}%
    \addcontentsline{toc}{chapter}{%
      \protect\texorpdfstring{\protect\MakeUppercase{\@chapapp\ \thechapter:\ #1}}%
                             {\@chapapp\ \thechapter: #1}}%
  \else
    \addcontentsline{toc}{chapter}{#1}%
  \fi
  \chaptermark{#1}%
  \addtocontents{lof}{\protect\addvspace{10\p@}}%
  \addtocontents{lot}{\protect\addvspace{10\p@}}%
  \@makechapterhead{#2}\@afterheading}
% danh mục hình/bảng: ghi đầy đủ "Hình 1.1:" / "Bảng 1.1:"
\renewcommand*\l@figure{\@dottedtocline{1}{0em}{5.6em}}
\renewcommand*\l@table{\@dottedtocline{1}{0em}{5.6em}}
\makeatother

% ---------------------------------------------------------------- tên tiếng Việt
\usepackage{etoolbox}
\gappto\captionsvietnamese{%
  \renewcommand\contentsname{Mục lục}%
  \renewcommand\listfigurename{Danh mục hình}%
  \renewcommand\listtablename{Danh mục bảng}%
  \renewcommand\bibname{Tài liệu tham khảo}%
  \renewcommand\figurename{Hình}%
  \renewcommand\tablename{Bảng}%
  \renewcommand\chaptername{Chương}%
  \renewcommand\appendixname{Phụ lục}%
}
\AtBeginDocument{%
  \renewcommand\lstlistingname{Đoạn mã}%
  \renewcommand\lstlistlistingname{Danh mục đoạn mã}}

% ---------------------------------------------------------------- tài liệu tham khảo
\usepackage[numbers, square, sort&compress]{natbib}
\setlength{\bibsep}{3pt plus 1pt}
\renewcommand{\bibsection}{%
  \chapter*{\bibname}\addcontentsline{toc}{chapter}{TÀI LIỆU THAM KHẢO}}

% ---------------------------------------------------------------- liên kết
\usepackage[hyphens]{url}
\usepackage[colorlinks=true, linkcolor=black, citecolor=black, urlcolor=black,
            bookmarksnumbered=true, unicode=true]{hyperref}
\usepackage{bookmark}

% ---------------------------------------------------------------- tiện ích
% tên tệp/hàm: chữ đơn cách; cho phép xuống dòng sau dấu "_" để không tràn lề
\newcommand{\codeus}{\textunderscore\allowbreak}
\newcommand{\code}[1]{\texttt{\small\let\_\codeus #1}}
\newcommand{\tn}[1]{\textit{#1}}                     % thuật ngữ tiếng Anh
\newcommand{\xor}{\oplus}
\newcommand{\hinh}[1]{Hình~\ref{#1}}
\newcommand{\bang}[1]{Bảng~\ref{#1}}
\newcommand{\doanma}[1]{Đoạn mã~\ref{#1}}
\newcommand{\pt}[1]{(\ref{#1})}
\newcommand{\muc}[1]{mục~\ref{#1}}
\newcommand{\ok}{\ding{51}}   % dấu tick (màu đen)
\newcommand{\no}{\ding{55}}   % dấu chéo (màu đen)
\usepackage{pifont}
% ô tiêu đề bảng có nền
\newcommand{\thead}[1]{\textbf{#1}}   % hàng tiêu đề bảng: in đậm, không tô nền
% khối ghi chú: một khung nền nhạt có tiêu đề in đậm
\newsavebox{\ghichubox}
\newenvironment{ghichu}[1][Nhận xét]{%
  \par\smallskip\noindent
  \begin{lrbox}{\ghichubox}%
  \begin{minipage}{\dimexpr\linewidth-12pt-2\fboxrule\relax}%
  \setlength{\parindent}{0pt}\setlength{\parskip}{3pt}%
  \textbf{\color{ctublue}#1.}\ }%
  {\end{minipage}\end{lrbox}%
  {\setlength{\fboxsep}{6pt}\fcolorbox{codeframe}{codebg}{\usebox{\ghichubox}}}%
  \par\smallskip}
